
% Adrianic memory controller candidate

\subsection{Separated Learning and Recall Timescales}

\begin{align}
\frac{\partial \chi}{\partial t}
&=
D_\chi\nabla^2\chi
+
L(t)\frac{\psi-\chi}{\tau_L},\\
\left.\frac{\partial\psi}{\partial t}\right|_{\rm recall}
&=
R(t)\mu_R(\chi-\psi).
\end{align}

Define phase-invariant overlap
\begin{align}
\mathcal O(a,b)
=
\frac{|\langle a,b\rangle|}
{\|a\|_2\|b\|_2+\epsilon}.
\end{align}

Define source-of-surprise signals
\begin{align}
S_\psi &= 1-\mathcal O(\psi_{t^-},\psi_t),\\
S_\chi &= 1-\mathcal O(\chi_{t^-},\chi_t),\\
E_{\psi\chi} &= 1-\mathcal O(\psi_t,\chi_t).
\end{align}

A transparent controller may use:

\begin{itemize}
\item If $S_\psi>S_\chi+\Delta$ and $E_{\psi\chi}>E_0$:
      set $L(t)=0$ temporarily and retain fast recall.
\item If $S_\chi>S_\psi+\Delta$ and $E_{\psi\chi}>E_0$:
      set $R(t)=0$ temporarily and re-teach $\chi$ from $\psi$.
\item If both $S_\psi$ and $S_\chi$ exceed a damage threshold:
      use symmetric/equal repair until a reliable source emerges.
\item Otherwise use a fast-recall/slow-learning normal mode.
\end{itemize}

For memory-source damage, a severity-scaled repair interval is
\begin{align}
T_{\rm repair}
=
\operatorname{clip}
\left[
\tau_L(0.5+2S_\chi),T_{\min},T_{\max}
\right].
\end{align}

In the current numerical benchmark, a useful operating regime was
$\mu_R=5$, $\tau_L=2$, with recall updates approximately 24 times more
frequent than learning updates during normal operation.
